Un dels exoplanetes amb més possibilitats d'acollir vida és el Ross 128~b. Gira al voltant de l'estrella Ross 128 amb un període orbital de 9,9 dies, en una òrbita pràcticament circular de radi $7,42\times 10^{6}\ \mathrm{km}$, i la seva massa és 1,35 vegades la massa de la Terra.

\figura[0.3]{fig1.png}

\begin{apartat}{1}
Calculeu la massa de l'estrella Ross 128.
\end{apartat}

\begin{apartat}{1}
Suposant que l'exoplaneta Ross 128~b tingui la mateixa densitat que la Terra, calculeu-ne el radi i el mòdul de la intensitat del camp gravitatori a la seva superfície.
\end{apartat}

\dades{%
  $G = 6,67\times 10^{-11}\ \mathrm{N\,m^2\,kg^{-2}}$.\\
  Massa de la Terra, $M_\mathrm{T} = 5,98\times 10^{24}\ \mathrm{kg}$.\\
  Radi de la Terra, $R_\mathrm{T} = 6,37\times 10^{6}\ \mathrm{m}$.}
